\chapter{Introduction}

An investor who has some capital and several assets to invest in (stocks, bonds,
exchange-traded funds, commodities) has to decide how much of the money goes into
each one. The set of assets with the proportion assigned to each is called a
portfolio. Prices of these assets change every day, so an allocation that was
reasonable yesterday may stop being so, and the investor has to decide again and
again. Managing a portfolio means making this sequence of decisions.

Consider a market with $m$ assets traded over $n$ periods, which in this work are
trading days. A portfolio is a vector $b=(b_1,\dots,b_m)$ where $b_i$ is the
fraction of wealth invested in asset $i$. In this project the fractions cannot be
negative (short positions are not allowed) and they add up to one, so all the
wealth is invested. The set of possible portfolios is
\[
\Delta_m=\Big\{\,b\in\mathbb{R}^m : b_i\ge 0,\ \textstyle\sum_{i=1}^m b_i=1\,\Big\}.
\]

The evolution of the market is described by the \emph{price relatives}. In period
$t$, component $i$ of the vector $x_t$ is the ratio between the closing price of
asset $i$ at $t$ and its price at $t-1$. A value of $1.02$ means the asset rose
2\,\% that day, and $0.97$ means it fell 3\,\%. If the portfolio held at the start
of the period is $b_t$, wealth is multiplied that day by $b_t\cdot x_t$, which is
the average of the $x_{t,i}$ weighted by the portfolio. Starting from a wealth of
one, after $n$ periods we have
\[
S_n=\prod_{t=1}^{n} b_t\cdot x_t .
\]

The Online Portfolio Selection (OLPS) problem consists in choosing $b_1,\dots,b_n$
so that $S_n$ is as large as possible. The word \emph{online} means that $b_t$ is
decided before $x_t$ is known, using only the prices of earlier periods
$x_1,\dots,x_{t-1}$, which is what happens in real life. An OLPS algorithm takes
the price history and returns the next day's portfolio.

This setting differs from the classical Markowitz approach \parencite{markowitz}, which
estimates the mean and covariance of returns and then optimises the portfolio using
those estimates. OLPS algorithms assume no statistical model of the market, and they
are judged directly by the wealth they would have achieved. The work is based on the
book \emph{Online Portfolio Selection: Principles and Algorithms} by \textcite{lihoi}, which gathers the algorithms of the field under a common
notation and is used as the reference throughout the project.
